\documentclass[12pt]{article}

\usepackage{sbc-template}
\usepackage{graphicx,url}
\usepackage[utf8]{inputenc}
\usepackage[brazil]{babel}
\usepackage{amssymb}
\usepackage{booktabs}
\usepackage{multirow}

\sloppy

\title{Análise de Sensibilidade ao Tamanho do Problema no Isolamento
do Efeito do \textit{Cache} L1 em GPU AMD Polaris}

\author{Henrique de Almeida Diniz\inst{1}, Pedro Henrique Ramos Costa\inst{1}}

\address{Instituto de Ciências Exatas e de Informática\\
  Pontifícia Universidade Católica de Minas Gerais (PUC Minas) -- Belo Horizonte, MG -- Brasil
  \email{henriquealdiniz@gmail.com, phramoscosta@gmail.com}
}

\begin{document}

\maketitle

\begin{abstract}
This report documents a problem-size sensitivity analysis of the cache-flag experiment on an
AMD Ellesmere GPU (Polaris/GCN, gfx803). Ten kernel variants, differing only in
the memory path taken by their loads, were measured at five matrix sizes from
$N=1024$ to $N=16384$, with 100 measured executions each, totalling 50 measured
points and about 12.6 hours of GPU time. All 50 points were numerically
validated. Three results stand out. The cost of the assembly construction is
nearly invariant with size (1.52$\times$ to 1.46$\times$), whereas the cost
attributed to the \texttt{glc} flag falls markedly with it (2.47$\times$ to
1.65$\times$), which shows that a single-size measurement does not characterise
the effect. The \texttt{slc} flag, negligible at small sizes (1.9\%), reaches
11.8\% at $N=16384$ --- behaviour consistent with an eviction policy rather than
a bypass, and detectable only once the working set greatly exceeds the L2. And
the dispersion between executions collapses from 4.0\% to 0.02\%, so the
bench's resolution improves by two orders of magnitude across the analysis.
\end{abstract}

\begin{resumo}
Este relatório documenta uma análise de sensibilidade ao tamanho do problema no experimento de
\textit{flags} de \textit{cache} em GPU AMD Ellesmere (Polaris/GCN, gfx803). Dez
variantes de \textit{kernel}, que diferem apenas no caminho de memória de seus
\textit{loads}, foram medidas em cinco tamanhos de matriz, de $N=1024$ a
$N=16384$, com 100 execuções medidas cada, totalizando 50 pontos e cerca de 12,6
horas de GPU. Os 50 pontos foram validados numericamente. Três resultados se
destacam. O custo da construção em \textit{assembly} é quase invariante com o
tamanho (1,52$\times$ a 1,46$\times$), ao passo que o custo atribuído à
\textit{flag} \texttt{glc} cai acentuadamente com ele (2,47$\times$ a
1,65$\times$), o que mostra que uma medida de tamanho único não caracteriza o
efeito. A \textit{flag} \texttt{slc}, desprezível em tamanhos pequenos (1,9\%),
alcança 11,8\% em $N=16384$ --- comportamento coerente com política de
\textit{eviction}, e não com \textit{bypass}, e detectável apenas quando o
conjunto de trabalho excede largamente o L2. E a dispersão entre execuções
colapsa de 4,0\% para 0,02\%, de modo que a resolução da bancada melhora duas
ordens de grandeza ao longo da análise.
\end{resumo}


\section{Objetivo}

Medições anteriores do custo de contornar o \textit{cache} L1 foram conduzidas
em um único tamanho de problema, $N=1024$. Duas razões tornam essa prática
insuficiente. A primeira é que $N=1024$ é múltiplo de 256 e, com armazenamento
\textit{row-major} de elementos de 4~bytes, produz \textit{stride} entre linhas
múltiplo de 1~KB, condição associada a \textit{conflict misses} no L1. A
segunda, e mais geral, é que nada garante que o custo de uma \textit{flag} de
\textit{cache} seja constante em relação ao tamanho do problema: à medida que o
conjunto de trabalho cresce, a fração dos acessos que o \textit{cache} conseguia
atender muda, e com ela o valor de contorná-lo.

Este relatório documenta a análise conduzida para responder a essa questão:
como cada fator isolado no desenho experimental varia quando o tamanho do
problema cresce por um fator de 4.096 em volume de dados.

\section{Desenho experimental}

Dez variantes do mesmo algoritmo --- multiplicação de matrizes \textit{naive},
um \textit{work-item} por elemento de saída --- foram medidas. Elas executam o
mesmo trabalho aritmético e diferem exclusivamente no caminho de memória de seus
\textit{loads} e na forma como esse caminho é solicitado ao compilador. Três
delas são \textbf{controles}, e existem para que os fatores confundidos possam
ser descontados:

\begin{itemize}
\item \textbf{Controle negativo} (variantes 2 e 3): dois \textit{kernels} de
fonte idêntico, que geram ISA idêntica instrução a instrução. A diferença
medida entre eles é integralmente artefato da bancada e fornece o \textbf{piso
de ruído} abaixo do qual nenhum efeito é afirmado.

\item \textbf{Controle da construção em \textit{assembly}} (variante 4): usa
\textit{inline assembly} sem \textit{flag} alguma. O bloco \texttt{\_\_asm\_\_}
é opaco ao otimizador, impede a redução de força no cálculo de endereço e
acrescenta trabalho de ALU por iteração. Comparar a variante 1 com a 4 mede esse
custo; comparar a 4 com a 5 mede o efeito da \textit{flag} já descontado dele.

\item \textbf{Controle da construção por intrinsic} (variante 9): usa
\texttt{llvm.amdgcn.raw.buffer.load.f32} com política de \textit{cache} nula.
Diferentemente do bloco de \textit{assembly}, o intrinsic tem semântica
conhecida pelo compilador, que pode escalonar em volta dele. Comparar a variante
9 com a 10 isola a \textit{flag} \texttt{glc} sozinha sobre código otimizado ---
combinação inalcançável pelas demais vias, já que \texttt{\_\_builtin\_nontemporal\_load}
emite \texttt{glc} e \texttt{slc} em conjunto, e o \textit{assembly} cobra a
barreira de otimização.
\end{itemize}

\subsection{Instrumentação}

Cada ponto consiste em 100 execuções medidas, precedidas de execuções de
aquecimento descartadas. A temporização usa \texttt{clGetEventProfilingInfo}
entre \texttt{COMMAND\_START} e \texttt{COMMAND\_END}, ou seja, mede
\textbf{apenas a execução do \textit{kernel}}: as transferências entre
\textit{host} e dispositivo ficam fora da janela cronometrada, por serem
idênticas entre as dez variantes e por estarem amortizadas em 100 execuções
sobre os mesmos \textit{buffers}. Os valores de GFLOPS deste relatório são,
portanto, de \textit{kernel}, e não fim a fim.

Registram-se média, mediana, desvio padrão, coeficiente de variação, percentil
95 e a deriva entre o primeiro e o último décimo das execuções, esta última para
detectar variação de \textit{clock} sob carga sustentada. As razões entre
cenários são calculadas sobre medianas, menos sensíveis a valores extremos.

\subsection{Validação}

Os 50 pontos foram verificados numericamente, em duas camadas. Cada variante é
comparada elemento a elemento com a saída da variante 1, em $O(N^2)$; e a
variante 1 é conferida contra uma referência calculada na CPU sobre uma amostra
aleatória de 4.096 posições, com semente fixa. A amostragem é necessária porque
a referência completa é $O(N^3)$ e, em $N=16384$, custaria cerca de três horas
de CPU --- mais do que a própria medição. Ela é suficiente porque as falhas em
questão não são pontuais: um descritor de recurso mal formado devolve zero em
todas as posições, e uma \textit{flag} mal emitida altera o caminho de memória
de todos os acessos.

A tolerância é proporcional a $N$, já que a soma acumula erro de ponto flutuante
ao longo de $N$ termos e um limite fixo produziria falso negativo em matrizes
grandes.

\subsection{Restrição de \textit{hardware}: o \textit{watchdog} de 60~s}

O \textit{driver} \texttt{amdgpu} declara travamento e reinicia a GPU quando um
trabalho de \textit{compute} ultrapassa 60~s. Em $N=16384$, um lançamento único
do \textit{kernel} mais lento foi estimado em cerca de 58~s, dentro da margem de
erro desse limite. Como a GPU utilizada é também a única saída de vídeo da
máquina, um reinício encerraria a sessão gráfica e a medição em curso.

A solução adotada foi fatiar a \textit{NDRange} em faixas de linhas por meio do
\texttt{global\_work\_offset}, com 2.048 linhas por lançamento em $N=16384$, o
que resulta em oito lançamentos de cerca de 7~s por execução. O trabalho total
permanece idêntico e os \textit{kernels} não requerem alteração, pois
\texttt{get\_global\_id(1)} já devolve o índice global com o deslocamento
incorporado. O custo acrescido é o de lançamento, medido em 0,02~ms por lote,
desprezível diante de segundos de execução. Antes de aplicar o modo em lotes ao
tamanho maior, verificou-se em $N=2048$ que ele produz resultado idêntico ao
lançamento único.

\subsection{Limite superior da análise}

O maior tamanho medido é $N=16384$. Tamanhos maiores não são alcançáveis nesta
máquina: em $N=32768$, as três matrizes exigiriam 12~GB de memória de vídeo,
contra os 8~GB disponíveis, e 16~GB de memória de sistema, contra 15~GB
instalados. Em $N=16384$ o experimento ocupa 3~GB de VRAM e 4~GB de RAM.

\section{Resultados}

A análise completa consumiu cerca de 12,6~horas de execução contínua, das
quais 11~horas no ponto $N=16384$. Os 50 pontos foram validados sem divergência.

\begin{table}[ht]
\centering
\caption{Mediana do tempo de execução, em milissegundos, por tamanho e cenário.
Cada valor resume 100 execuções medidas.}
\label{tab:medianas}
\small
\begin{tabular}{@{}llrrrrr@{}}
\toprule
\textbf{Cenário} & \textbf{\textit{Flags}} & \textbf{1024} & \textbf{2048} &
\textbf{4096} & \textbf{8192} & \textbf{16384} \\
\midrule
1. normal (C puro) & GLC=0 SLC=0 & 6.99 & 55.84 & 449.0 & 3693 & 29934 \\
2. \texttt{volatile} & GLC=1$^\ast$ & 13.40 & 92.12 & 744.5 & 5940 & 47671 \\
3. \texttt{volatile} (idem 2) & GLC=1$^\ast$ & 13.28 & 92.39 & 745.3 & 5952 & 47714 \\
4. controle \textit{asm} & GLC=0 SLC=0 & 10.62 & 84.37 & 675.5 & 5429 & 43642 \\
5. \textit{asm} \texttt{glc} & GLC=1 & 12.73 & 92.31 & 719.0 & 5747 & 46343 \\
6. \textit{asm} \texttt{glc slc} & GLC=1 SLC=1 & 12.98 & 93.46 & 737.3 & 6048 & 51816 \\
7. \texttt{nontemporal} & GLC=1 SLC=1 & 12.42 & 86.62 & 666.8 & 5397 & 47536 \\
8. \textit{tiled} (LDS) & GLC=0 SLC=0 & 2.11 & 16.24 & 129.5 & 1078 & 9369 \\
9. controle intrinsic & GLC=0 SLC=0 & 5.42 & 41.19 & 349.8 & 2896 & 23630 \\
10. intrinsic \texttt{glc} & GLC=1 & 13.37 & 85.71 & 623.2 & 4895 & 39014 \\

\bottomrule
\end{tabular}

\vspace{2pt}
{\footnotesize $^\ast$ depende do compilador: \texttt{glc} sob o Clang/LLVM do
ROCm, nenhuma \textit{flag} sob o AMDOCL do Windows.}
\end{table}

\subsection{A dispersão colapsa com o tamanho}

A Tabela~\ref{tab:cv} mostra o coeficiente de variação entre as 100 execuções de
cada ponto. Ele cai de cerca de 4\% em $N=1024$ para 0,02\% em $N=16384$: com
\textit{kernels} de dezenas de segundos, o ruído de lançamento e de escalonamento
torna-se irrelevante diante da duração medida.

\begin{table}[ht]
\centering
\caption{Coeficiente de variação entre as 100 execuções, em porcentagem.}
\label{tab:cv}
\small
\begin{tabular}{@{}lrrrrr@{}}
\toprule
\textbf{Cenário} & \textbf{1024} & \textbf{2048} & \textbf{4096} &
\textbf{8192} & \textbf{16384} \\
\midrule
1. normal (OpenCL C) & 3.99 & 0.62 & 0.28 & 0.03 & 0.02 \\
2. \texttt{volatile} & 3.02 & 0.84 & 0.17 & 0.07 & 0.03 \\
3. \texttt{volatile} (idem 2) & 2.66 & 1.27 & 0.14 & 0.14 & 0.03 \\
4. controle \textit{asm} & 1.70 & 0.21 & 0.02 & 0.01 & 0.06 \\
5. \textit{asm} \texttt{glc} & 3.09 & 1.40 & 0.08 & 0.04 & 0.03 \\
6. \textit{asm} \texttt{glc slc} & 2.66 & 1.65 & 0.46 & 0.04 & 0.07 \\
7. \texttt{nontemporal} & 3.24 & 0.77 & 0.15 & 0.05 & 0.20 \\
8. \textit{tiled} (LDS) & 2.65 & 3.29 & 0.69 & 0.57 & 0.47 \\
9. controle intrinsic & 2.63 & 1.30 & 0.14 & 0.05 & 0.03 \\
10. intrinsic \texttt{glc} & 5.31 & 2.33 & 0.37 & 0.08 & 0.05 \\

\bottomrule
\end{tabular}
\end{table}

A consequência é metodológica: a resolução da bancada melhora duas ordens de
grandeza ao longo da análise. Efeitos de poucos por cento, que em $N=1024$ se
confundem com o ruído, tornam-se inequívocos em $N=16384$. O piso de ruído
medido pelo par de \textit{kernels} idênticos acompanha esse comportamento,
mantendo-se em torno de 0,1\% a partir de $N=2048$ (Tabela~\ref{tab:razoes}).

\subsection{Decomposição dos fatores}

\begin{table}[ht]
\centering
\caption{Razões entre cenários, calculadas sobre as medianas. Valores acima de
1 indicam que o cenário do numerador é mais lento.}
\label{tab:razoes}
\small
\begin{tabular}{@{}lrrrrr@{}}
\toprule
\textbf{Fator isolado} & \textbf{1024} & \textbf{2048} & \textbf{4096} &
\textbf{8192} & \textbf{16384} \\
\midrule
Piso de ruído (3/2): \textit{kernels} de ISA idêntica & 0.991 & 1.003 & 1.001 & 1.002 & 1.001 \\
Custo da construção em \textit{asm} (4/1) & 1.520 & 1.511 & 1.504 & 1.470 & 1.458 \\
Custo da construção por intrinsic (9/1) & 0.776 & 0.738 & 0.779 & 0.784 & 0.789 \\
Efeito de \texttt{glc}, base em \textit{asm} (5/4) & 1.199 & 1.094 & 1.064 & 1.059 & 1.062 \\
Efeito de \texttt{glc}, base em intrinsic (10/9) & 2.466 & 2.081 & 1.781 & 1.691 & 1.651 \\
Efeito de \texttt{glc slc} sem \textit{asm} (7/1) & 1.776 & 1.551 & 1.485 & 1.461 & 1.588 \\
Efeito de \texttt{slc} isolado (6/5) & 1.019 & 1.012 & 1.026 & 1.052 & 1.118 \\
\texttt{volatile} sob ROCm (2/1) & 1.917 & 1.650 & 1.658 & 1.608 & 1.593 \\
LDS contra memória global (1/8) & 3.305 & 3.438 & 3.466 & 3.425 & 3.195 \\

\bottomrule
\end{tabular}
\end{table}

A Tabela~\ref{tab:razoes} separa fatores que se comportam de maneiras
qualitativamente distintas.

\textbf{O custo da construção é estável.} A barreira de otimização imposta pelo
bloco de \textit{assembly} custa entre 1,52$\times$ e 1,46$\times$ em toda a
faixa, com tendência suave de queda. O intrinsic, por sua vez, é
\textbf{mais rápido} que o código em C puro por um fator praticamente constante
de 0,78$\times$ --- resultado atribuível ao formato MUBUF, que mantém o
descritor em registradores escalares e usa deslocamento de 32~bits por
\textit{lane}, reduzindo a alocação de 11 para 7 registradores vetoriais e
elevando a ocupância. Que ambos os fatores sejam quase invariantes é o que
permite descontá-los com segurança.

\textbf{O custo da \textit{flag} \texttt{glc} não é estável.} Medido sobre a
base em intrinsic, cai de 2,47$\times$ em $N=1024$ para 1,65$\times$ em
$N=16384$; sobre a base em \textit{assembly}, de 1,20$\times$ para 1,06$\times$.
A interpretação é direta: à medida que o conjunto de trabalho cresce, a fração
dos acessos que o L1 conseguia atender diminui, e portanto diminui também o que
se perde ao contorná-lo. \textbf{Uma medida em um único tamanho não caracteriza
o efeito do \textit{cache}} --- e, tomada isoladamente, o superestima em relação
ao regime assintótico.

\textbf{As duas bases discordam por construção.} O mesmo bit medido contra bases
diferentes produz 1,65$\times$ e 1,06$\times$ em $N=16384$. Não há contradição:
a base em \textit{assembly} já está cerca de 1,46$\times$ mais lenta, e uma base
mais lenta dilui a penalidade relativa do \textit{cache}. As duas razões
respondem a perguntas distintas --- quanto custa perder o L1 em código que se
pretende efetivamente usar, e qual o efeito do bit isolado de todo o resto.
Relatar uma delas sem declarar a base torna o número ininterpretável.

\subsection{O efeito de \texttt{slc} aparece apenas em conjuntos grandes}

O resultado mais informativo da análise está na linha do \texttt{slc} isolado.
A \textit{flag} custa 1,9\% em $N=1024$, valor próximo ao piso de ruído e que em
medições anteriores oscilou de veredito entre execuções independentes. Com o
crescimento do problema, o efeito cresce de forma monotônica e alcança
\textbf{11,8\% em $N=16384$}, duas ordens de grandeza acima do piso de ruído
daquele ponto.

Esse comportamento é exatamente o previsto para uma \textbf{política de
substituição}, e não para um \textit{bypass}. Marcar uma linha como candidata
preferencial a \textit{eviction} não produz efeito algum quando o conjunto de
trabalho cabe no \textit{cache}, pois não há disputa por espaço; o efeito só se
manifesta quando há mais dados concorrendo do que capacidade disponível. Em
$N=16384$, as três matrizes somam 3~GB disputando 2~MB de L2.

O achado corrobora, por via independente e sob carga real, a conclusão obtida em
experimento dedicado de vazão: \texttt{SLC=1} não realiza \textit{bypass} do L2.
Um \textit{bypass} produziria penalidade grande e aproximadamente constante,
presente já nos tamanhos pequenos; o que se observa é o oposto --- efeito nulo
enquanto há folga no \textit{cache} e crescente conforme a folga desaparece.

\subsection{Demais observações}

O ganho da variante com memória local (LDS) sobre o acesso à memória global
mantém-se entre 3,2$\times$ e 3,5$\times$ em toda a faixa, com leve queda no
ponto extremo. O desempenho absoluto do \textit{kernel} de referência decai
suavemente com o tamanho, de 307 para 294 GFLOPS, ao passo que o da variante com
intrinsic e \texttt{glc} \textbf{sobe} de 161 para 226 GFLOPS: as duas curvas
convergem, o que é a mesma tendência de queda do custo da \textit{flag} vista de
outro ângulo.

O \texttt{volatile} sob o ROCm, que emite \texttt{glc} e adicionalmente
serializa os \textit{loads}, acompanha a queda geral, de 1,92$\times$ para
1,59$\times$.

\section{Limitações}

\textbf{Todos os tamanhos medidos são múltiplos de 256.} A série dobra $N$ a
cada ponto a partir de 1024, de modo que a condição associada a \textit{conflict
misses} no L1 --- \textit{stride} entre linhas múltiplo de 1~KB --- vale em todos
os cinco pontos. As tendências relatadas são, portanto, tendências \emph{dentro}
desse regime. Medições em tamanhos fora da grade de 256 apresentaram efeito de
\texttt{glc} sensivelmente menor, o que indica que a comparação entre regimes
exige série própria, com passo fino e tamanhos não alinhados.

\textbf{Cem execuções por ponto.} O número foi escolhido para tornar a análise
viável: mil execuções, usadas em medições de tamanho único, custariam mais de
cinco dias apenas em $N=16384$. Em tamanhos pequenos isso amplia o intervalo de
confiança em relação às medições anteriores; em tamanhos grandes é irrelevante,
dada a dispersão de 0,02\%.

\textbf{Uma única carga e uma única geração de \textit{hardware}.} Todos os
pontos vêm da multiplicação de matrizes \textit{naive} em uma GPU gfx803. A
generalização das tendências para outros padrões de acesso e outras arquiteturas
não está estabelecida por estes dados.

\textbf{Apenas leituras.} O comportamento das \textit{flags} em instruções de
escrita permanece em aberto, questão relevante por o L1 do GFX8 ser
\textit{write-through}.

\section{Reprodução}

O programa aceita o número de execuções, o tamanho da matriz, o número de
execuções de aquecimento e o número de linhas por lote:

\begin{verbatim}
source setup/env_rocm.sh
g++ -Wall -Wextra -O2 experimentos/matmul.cpp -o output/matmul -lOpenCL

./output/matmul 100  1024 5         # lancamento unico
./output/matmul 100 16384 3 2048    # lotes de 2048 linhas
\end{verbatim}

O quarto argumento é obrigatório a partir de $N=16384$, pelas razões expostas na
Seção~2.4. Passar zero, ou omiti-lo, faz o programa usar um lançamento único.
Passar zero como primeiro argumento ativa o modo adaptativo, no qual o número de
execuções é derivado de um orçamento de tempo por cenário.

A análise completa é reproduzida por \texttt{tools/cadeia\_sweep.sh}, que
encadeia os tamanhos, verifica a integridade do que já foi gravado e valida o
modo em lotes antes de prosseguir para o tamanho maior. Os resultados são
acumulados em \texttt{resultados/matmul\_sweep.csv}, com uma linha por par
(tamanho, cenário) e gravação imediata, de modo que uma interrupção não perde o
que já foi medido. As amostras individuais de cada tamanho ficam em
\texttt{resultados/matmul\_amostras\_N$\langle$N$\rangle$.csv}, para análise de
distribuição.

\end{document}
